% ============================================================================
% APPENDIX: STORY ARCS & LOOP PATTERNS — Visual Reference Library
% ============================================================================

\section{Story Arcs \& Loop Patterns: Visual Reference Library}

\begin{center}
{\small\color{slate}\itshape The complete collection of narrative structures and argumentative patterns. Use these as building blocks for your presentations.}
\end{center}

% ============================================================================
\subsection{Story Arcs}

Story arcs define the macro-structure of your presentation --- the emotional journey from opening to close. Choose the arc that matches your intent and audience.

\vspace{8pt}

\sourcefig[0.80]{infographics/arc-sparkline.jpg}{The Sparkline (Nancy Duarte): alternating ``What Is'' and ``What Could Be.'' Best for keynotes, vision speeches, and all-hands meetings.}

\sourcefig[0.80]{infographics/arc-heros-journey.jpg}{The Hero's Journey (Joseph Campbell): the audience is the hero, you are the mentor. Best for brand narratives, client relationships, and transformation stories.}

\sourcefig[0.80]{infographics/arc-consultants-gambit.jpg}{The Consultant's Gambit: structured problem-solving narrative. Best for strategy presentations, board decks, and analytical recommendations.}

\sourcefig[0.80]{infographics/arc-problem-agitate.jpg}{Problem-Agitate-Solution: intensify the pain before offering relief. Best for sales pitches, change management, and budget requests.}

\sourcefig[0.80]{infographics/arc-sequoia-pitch.jpg}{The Sequoia Pitch: the framework behind the world's most funded startups. Problem, Solution, Market, Business Model, Ask.}

\sourcefig[0.80]{infographics/arc-golden-circle.jpg}{The Golden Circle (Simon Sinek): Why before How before What. Best for purpose-driven communication and brand positioning.}

\clearpage

% ============================================================================
\subsection{Loop Patterns}

Loops are the argumentative building blocks within each story arc --- sequences of 3--10 slides that form one complete argument. Master these five core patterns and you can construct any presentation.

\vspace{8pt}

\begin{center}
\begin{minipage}[t]{0.45\textwidth}
\sourcefig[1.0]{infographics/loop-logic-chain.jpg}{The Logic Chain: deductive reasoning. Premise + Premise = Inevitable Conclusion. For sceptical audiences who need step-by-step proof.}
\end{minipage}
\hfill
\begin{minipage}[t]{0.45\textwidth}
\sourcefig[1.0]{infographics/loop-pattern-hunter.jpg}{The Pattern Hunter: inductive reasoning. Three data points reveal one pattern. For data-driven arguments.}
\end{minipage}
\end{center}

\vspace{8pt}

\begin{center}
\begin{minipage}[t]{0.45\textwidth}
\sourcefig[1.0]{infographics/loop-aha-moment.jpg}{The Aha Moment: build context, then deliver the reveal. For insights that need setup to land with impact.}
\end{minipage}
\hfill
\begin{minipage}[t]{0.45\textwidth}
\sourcefig[1.0]{infographics/loop-before-after.jpg}{Before/After: show the transformation. Grey world becomes teal. For change management and solution proposals.}
\end{minipage}
\end{center}

\vspace{8pt}

\begin{center}
\begin{minipage}[t]{0.45\textwidth}
\sourcefig[1.0]{infographics/loop-so-what.jpg}{The So What? Cascade: drill three levels deep to find the real insight beneath the surface data.}
\end{minipage}
\end{center}

\vspace{12pt}

\begin{insightbox}[The Full Collection]
This appendix shows 6 of the 20 story arcs and 5 of the 54 loop patterns available in the Storymakers Playbook. The complete visual library --- including patterns like The Iceberg, The Trojan Horse, The Pre-Mortem, The Three Horizons, and 45 more --- is available at \textbf{storymakersbyanlak.com}.
\end{insightbox}

\clearpage
